As the example above indicates, we will need to assume that the boundary has some smoothness properties in order to get a boundary regularity result (also called a global regularity result). In order to emphasize the most important techniques in the proof (breaking up the domain using a partition of unity and mapping the pieces containing portions of the boundary to a half-space) we will give the proof only for second-order scalar equations and in the proof we will ignore lower-order terms.

\begin{Theorem}{Global regularity}
\label{thm:global-boundary-regularity}
Suppose that the hypotheses of Theorem~\ref{thm:second-order-interior-regularity} hold and that in addition $\partial\Omega$ is of class $C^2$. Then $u\in H^2(\Omega)$ and
\begin{equation}\label{eq:global-boundary-regularity-h2-estimate}
\|u\|_{H^2(\Omega)}\leq C(\|u\|_{L^2(\Omega)}+\|f\|_{L^2(\Omega)}).
\end{equation}

\end{Theorem}

The proof of this result is rather long and involved, so we will break it up by proving a number of preliminary lemmas. One of our basic techniques is to decompose the domain into pieces using a partition of unity and ``flattening out\textquotesingle\textquotesingle\ any portion of the boundary. As we see in our first lemma (which is essentially a version of the main result in the case where the boundary is already flat) a flat boundary allows us to use difference quotients to our advantage.

\begin{Lemma}
\label{lem:flat-boundary-local-regularity}
Let $R>0$, $\lambda\in(0,1)$, and define
\begin{equation}\label{eq:flat-boundary-outer-half-ball}
D^+:=B_R(0)\cap\{\bm{x}\in\mathbb{R}^n\mid x_n>0\},
\end{equation}

\begin{equation}\label{eq:flat-boundary-inner-half-ball}
Q^+:=B_{\lambda R}(0)\cap\{\bm{x}\in\mathbb{R}^n\mid x_n>0\},
\end{equation}

\begin{equation}\label{eq:flat-boundary-source-isolated-empty-number}\end{equation}
Let $L$ be a uniformly elliptic second-order operator of the form\[L(\bm{x},D)u:=\sum_{i,j=1}^n\frac{\partial}{\partial x_i}a_{ij}(\bm{x})\frac{\partial u}{\partial x_j}+\sum_{i=1}^nb_i(\bm{x})\frac{\partial u}{\partial x_i}+c(\bm{x})u\]with corresponding bilinear form\[\begin{aligned}
B[v,u]&:=-\sum_{i,j=1}^n\int_{D^+}a_{ij}(\bm{x})u_{x_j}v_{x_i}\,\mathrm{d}\bm{x}+\sum_{i=1}^n\int_{D^+}b_i(\bm{x})u_{x_i}v\,\mathrm{d}\bm{x}\\
&\quad+\int_{D^+}c(\bm{x})uv\,\mathrm{d}\bm{x}.
\end{aligned}\]Suppose the coefficients satisfy $a_{ij}\in W^{1,\infty}(D^+)$, $b_i,c\in L^\infty(D^+)$ and that $f\in L^2(D^+)$. Suppose $u\in H^1(D^+)$ satisfies the variational equation
\begin{equation}\label{eq:flat-boundary-variational-equation}
B[v,u]=(v,f)
\end{equation}
for all $v\in H_0^1(D^+)$ and that $u\equiv0$ in the sense of trace on $\{\bm{x}\in\mathbb{R}^n\mid x_n=0\}$. Then $u\in H^2(Q^+)$ and there exists a constant $C$ depending on $R$ such that
\begin{equation}\label{eq:flat-boundary-local-h2-estimate}
\|u\|_{H^2(Q^+)}\leq C(\|f\|_{L^2(D^+)}+\|u\|_{L^2(D^+)}).
\end{equation}

\end{Lemma}
